\documentclass[11pt]{article}

%------------------------------------------------
% Packages
%------------------------------------------------
\usepackage[a4paper,margin=2.5cm]{geometry}
\usepackage{amsmath,amssymb,amsthm}
\usepackage{mathtools}
\usepackage{enumitem}
\usepackage{hyperref}

\hypersetup{
    colorlinks=true,
    linkcolor=blue,
    citecolor=blue,
    urlcolor=blue
}

%------------------------------------------------
% Title Information
%------------------------------------------------

\title{\bfseries Adaptive Block-Coordinate Forward-Backward: Particular case of multilevel block strategy for Image Reconstruction}

\author{Nelly Pustelnik}
%thanks{Department of Mathematics, University A,
%\texttt{first.author@university.edu}}
%\and
%Second Author\thanks{Institute of Applied Mathematics, University B,
%\texttt{second.author@university.edu}}
%}

\date{}

%------------------------------------------------

\begin{document}

\maketitle

\vspace{-1em}

\noindent
\textbf{Keywords:}
Adaptive Block-Coordinate Forward-Backward: Particular case of multilevel block strategy for Image Reconstruction

\vspace{0.5em}


\vspace{1em}

%%%%%%%%%%%%%%%%%%%%%%%%%%%%%%%%%%%%%%%%%%%%%%%%%%
\section*{Abstract}
Classical first-order optimization methods for imaging inverse problems struggle to scale efficiently with increasing image resolution. Here, we introduce an adaptive multiresolution block coordinate Forward-Backward algorithm for image restoration. The selection of multiresolution blocks is guided by the local magnitude of the proximal update, using a stochastic, non-smooth Gauss-Southwell rule applied to the image wavelet decomposition. This adaptive approach dynamically balances updates across scales, prioritizing coarse or fine blocks based on the degradation characteristics. As a result, the proposed method automatically adjusts to varying levels of blur and noise without requiring a predefined hierarchical update scheme. 
%%%%%%%%%%%%%%%%%%%%%%%%%%%%%%%%%%%%%%%%%%%%%%%%%%
\vspace{1em}

\noindent
%\textbf{Acknowledgements.}
%This work was supported by ...

%%%%%%%%%%%%%%%%%%%%%%%%%%%%%%%%%%%%%%%%%%%%%%%%%%
\vspace{1em}

%\begin{thebibliography}{9}


%\end{thebibliography}

\end{document}